% 附录 F — 第三方软件许可（精译重写）
\biappendix{第三方软件许可}{Third-Party Software Licenses}
\label{chap:lic}

\section{引言}
\subsection*{Introduction}

RedHawk 产品包含需保留下列版权声明与许可条款的第三方软件。
本译本对各组件给出\textbf{中文导读}（许可类型、使用义务、与专有代码组合注意点），
长篇法律正文保留英文原文于 \texttt{lstlisting} 环境中，便于对照查阅。

\section{第三方组件概览}
\subsection*{Third-Party Components Overview}

\begin{longtable}{@{}p{0.30\textwidth}p{0.18\textwidth}p{0.46\textwidth}@{}}
\toprule
\textbf{组件/许可} & \textbf{类型} & \textbf{中文导读} \\
\midrule
\endhead
Common Public License (CPL) 1.0 & CPL & IBM 发起的开源许可证；允许使用、修改与再分发，但须保留版权声明与许可条款；商业分发者对最终用户承担额外担保义务。 \\
MIT License — ZHU Kaidi (2017) & MIT & 宽松 MIT 许可：可自由使用、修改、合并与再分发，须保留版权与许可声明；软件按“原样”提供，不附带担保。 \\
Boost Software License & Boost & 允许自由使用与再分发，须保留版权声明；免责声明排除担保责任。 \\
Jorn Lind-Nielsen 组件 & Custom & 1996–2002 版权；按作者许可条款使用与再分发。 \\
GAlib — Matthew Wall & Custom & 遗传算法库；非公有领域，按作者/MIT 条款使用。 \\
GIFLIB — Eric S. Raymond & MIT-like & GIF 图像库；MIT 风格宽松许可。 \\
LAPACK / BLAS & BSD & 田纳西大学、伯克利等机构数值线性代数库；BSD 风格许可。 \\
libarchive — Tim Kientzle & BSD & 归档/压缩文件读写库；BSD 风格许可。 \\
EFL / Evas & BSD & 嵌入式图形库；BSD 风格许可。 \\
Tcl/Tk & BSD & 脚本语言与 GUI 工具包；BSD 风格许可。 \\
METIS & Apache-2.0 & 图划分库；Apache License 2.0。 \\
Independent JPEG Group (IJG) & IJG & JPEG 编解码参考实现；须遵守 IJG 许可声明。 \\
Mozilla Public License 2.0 & MPL-2.0 & 部分 GUI/库组件；修改文件须以 MPL 发布，可与专有代码文件级隔离组合。 \\
OpenSSL / SSLeay & OpenSSL & 加密与 SSL/TLS 库；双重许可，须保留版权与免责。 \\
Mersenne Twister & BSD & 伪随机数生成器；BSD 风格许可。 \\
UMFPACK / SuiteSparse & LGPL & 稀疏矩阵求解；部分组件为 LGPL。 \\
Open MPI & BSD & MPI 实现；BSD 风格再分发条款。 \\
CSPARSE & LGPL & 稀疏矩阵库；GNU LGPL。 \\
SuperLU & BSD & UC Berkeley 等机构数值库；BSD 风格许可。 \\
zlib & zlib & 数据压缩库；zlib 宽松许可。 \\
\bottomrule
\end{longtable}

\section{许可正文（英文原文）}
\subsection*{License Texts (English)}

\begin{noteBox}
以下为原书附录 F 收录的第三方许可正文。完整法律文本以 ANSYS 官方发布包为准。
\end{noteBox}

\subsection{THE ACCOMPANYING PROGRAM IS PROVIDED UNDER THE TERMS OF THIS}

\begin{lstlisting}[basicstyle=\ttfamily\scriptsize,breaklines=true]
THE ACCOMPANYING PROGRAM IS PROVIDED UNDER THE TERMS OF THIS
              COMMON PUBLIC LICENSE ("AGREEMENT"). ANY USE, REPRODUCTION OR
              DISTRIBUTION OF THE PROGRAM CONSTITUTES RECIPIENT'S ACCEPTANCE OF
              THIS AGREEMENT.
              1. DEFINITIONS
              "Contribution" means:
              a) in the case of the initial Contributor, the initial code and documentation distributed
              under this Agreement, and
              b) in the case of each subsequent Contributor:
              i) changes to the Program, and
              ii) additions to the Program;
              where such changes and/or additions to the Program originate from and are distributed by
              that particular Contributor. A Contribution 'originates' from a Contributor if it was added to
              the Program by such Contributor itself or anyone acting on such Contributor's behalf.
              Contributions do not include additions to the Program which: (i) are separate modules of
              software distributed in conjunction with the Program under their own license agreement,
              and (ii) are not derivative works of the Program.
              "Contributor" means any person or entity that distributes the Program.
              "Licensed Patents " mean patent claims licensable by a Contributor which are necessarily
              infringed by the use or sale of its Contribution alone or when combined with the Program.
              "Program" means the Contributions distributed in accordance with this Agreement.
              "Recipient" means anyone who receives the Program under this Agreement, including all
              Contributors.
              2. GRANT OF RIGHTS
              a) Subject to the terms of this Agreement, each Contributor hereby grants Recipient a
              non-exclusive, worldwide, royalty-free copyright license to reproduce, prepare derivative
              works of, publicly display, publicly perform, distribute and sublicense the Contribution of
              such Contributor, if any, and such derivative works, in source code and object code form.
              b) Subject to the terms of this Agreement, each Contributor hereby grants Recipient a
              non-exclusive, worldwide, royalty-free patent license under Licensed Patents to make,
              use, sell, offer to sell, import and otherwise transfer the Contribution of such Contributor, if
              any, in source code and object code form. This patent license shall apply to the
              combination of the Contribution and the Program if, at the time the Contribution is added
              by the Contributor, such addition of the Contribution causes such combination to be
              covered by the Licensed Patents. The patent license shall not apply to any other
              combinations which include the Contribution. No hardware per se is licensed hereunder.
              c) Recipient understands that although each Contributor grants the licenses to its
              Contributions set forth herein, no assurances are provided by any Contributor that the
              Program does not infringe the patent or other intellectual property rights of any other
APPENDIX F — Third-Party Software Licenses                                      RedHawk User Manual | 962
entity. Each Contributor disclaims any liability to Recipient for claims brought by any other
              entity based on infringement of intellectual property rights or otherwise. As a condition to
              exercising the rights and licenses granted hereunder, each Recipient hereby assumes
              sole responsibility to secure any other intellectual property rights needed, if any. For
              example, if a third party patent license is required to allow Recipient to distribute the
              Program, it is Recipient's responsibility to acquire that license before distributing the
              Program.
              d) Each Contributor represents that to its knowledge it has sufficient copyright rights in its
              Contribution, if any, to grant the copyright license set forth in this Agreement.
              3. REQUIREMENTS
              A Contributor may choose to distribute the Program in object code form under its own
              license agreement, provided that:
              a) it complies with the terms and conditions of this Agreement; and
              b) its license agreement:
              i) effectively disclaims on behalf of all Contributors all warranties and conditions, express
              and implied, including warranties or conditions of title and non-infringement, and implied
              warranties or conditions of merchantability and fitness for a particular purpose;
              ii) effectively excludes on behalf of all Contributors all liability for damages, including
              direct, indirect, special, incidental and consequential damages, such as lost profits;
              iii) states that any provisions which differ from this Agreement are offered by that
              Contributor alone and not by any other party; and
              iv) states that source code for the Program is available from such Contributor, and informs
              licensees how to obtain it in a reasonable manner on or through a medium customarily
              used for software exchange.
              When the Program is made available in source code form:
              a) it must be made available under this Agreement; and
              b) a copy of this Agreement must be included with each copy of the Program.
              Contributors may not remove or alter any copyright notices contained within the Program.
              Each Contributor must identify itself as the originator of its Contribution, if any, in a
              manner that reasonably allows subsequent Recipients to identify the originator of the
              Contribution.
              4. COMMERCIAL DISTRIBUTION
              Commercial distributors of software may accept certain responsibilities with respect to
              end users, business partners and the like. While this license is intended to facilitate the
              commercial use of the Program, the Contributor who includes the Program in a
              commercial product offering should do so in a manner which does not create potential
              liability for other Contributors. Therefore, if a Contributor includes the Program in a
              commercial product offering, such Contributor ("Commercial Contributor") hereby agrees
              to defend and indemnify every other Contributor ("Indemnified Contributor") against any
              losses, damages and costs (collectively "Losses") arising from claims, lawsuits and other
              legal actions brought by a third party against the Indemnified Contributor to the extent
              caused by the acts or omissions of such Commercial Contributor in connection with its
              distribution of the Program in a commercial product offering. The obligations in this
              section do not apply to any claims or Losses relating to any actual or alleged intellectual
              property infringement. In order to qualify, an Indemnified Contributor must: a) promptly
              notify the Commercial Contributor in writing of such claim, and b) allow the Commercial
              Contributor to control, and cooperate with the Commercial Contributor in, the defense and
              any related settlement negotiations. The Indemnified Contributor may participate in any
              such claim at its own expense.
ANSYS, Inc.

APPENDIX F — Third-Party Software Licenses                                    RedHawk User Manual | 963
For example, a Contributor might include the Program in a commercial product offering,
              Product X. That Contributor is then a Commercial Contributor. If that Commercial
              Contributor then makes performance claims, or offers warranties related to Product X,
              those performance claims and warranties are such Commercial Contributor's
              responsibility alone. Under this section, the Commercial Contributor would have to defend
              claims against the other Contributors related to those performance claims and warranties,
              and if a court requires any other Contributor to pay any damages as a result, the
              Commercial Contributor must pay those damages.
              5. NO WARRANTY
              EXCEPT AS EXPRESSLY SET FORTH IN THIS AGREEMENT, THE PROGRAM IS
              PROVIDED ON AN "AS IS" BASIS, WITHOUT WARRANTIES OR CONDITIONS OF
              ANY KIND, EITHER EXPRESS OR IMPLIED INCLUDING, WITHOUT LIMITATION, ANY
              WARRANTIES OR CONDITIONS OF TITLE, NON-INFRINGEMENT,
              MERCHANTABILITY OR FITNESS FOR A PARTICULAR PURPOSE. Each Recipient is
              solely responsible for determining the appropriateness of using and distributing the
              Program and assumes all risks associated with its exercise of rights under this
              Agreement, including but not limited to the risks and costs of program errors, compliance
              with applicable laws, damage to or loss of data, programs or equipment, and
              unavailability or interruption of operations.
              6. DISCLAIMER OF LIABILITY
              EXCEPT AS EXPRESSLY SET FORTH IN THIS AGREEMENT, NEITHER RECIPIENT
              NOR ANY CONTRIBUTORS SHALL HAVE ANY LIABILITY FOR ANY DIRECT,
              INDIRECT, INCIDENTAL, SPECIAL, EXEMPLARY, OR CONSEQUENTIAL DAMAGES
              (INCLUDING WITHOUT LIMITATION LOST PROFITS), HOWEVER CAUSED AND ON
              ANY THEORY OF LIABILITY, WHETHER IN CONTRACT, STRICT LIABILITY, OR TORT
              (INCLUDING NEGLIGENCE OR OTHERWISE) ARISING IN ANY WAY OUT OF THE
              USE OR DISTRIBUTION OF THE PROGRAM OR THE EXERCISE OF ANY RIGHTS
              GRANTED HEREUNDER, EVEN IF ADVISED OF THE POSSIBILITY OF SUCH
              DAMAGES.
              7. GENERAL
              If any provision of this Agreement is invalid or unenforceable under applicable law, it shall
              not affect the validity or enforceability of the remainder of the terms of this Agreement,
              and without further action by the parties hereto, such provision shall be reformed to the
              minimum extent necessary to make such provision valid and enforceable.
              If Recipient institutes patent litigation against a Contributor with respect to a patent
              applicable to software (including a cross-claim or counterclaim in a lawsuit), then any
              patent licenses granted by that Contributor to such Recipient under this Agreement shall
              terminate as of the date such litigation is filed. In addition, if Recipient institutes patent
              litigation against any entity (including a cross-claim or counterclaim in a lawsuit) alleging
              that the Program itself (excluding combinations of the Program with other software or
              hardware) infringes such Recipient's patent(s), then such Recipient's rights granted under
              Section 2(b) shall terminate as of the date such litigation is filed.
              All Recipient's rights under this Agreement shall terminate if it fails to comply with any of
              the material terms or conditions of this Agreement and does not cure such failure in a
              reasonable period of time after becoming aware of such noncompliance. If all Recipient's
              rights under this Agreement terminate, Recipient agrees to cease use and distribution of
              the Program as soon as reasonably practicable. However, Recipient's obligations under
              this Agreement and any licenses granted by Recipient relating to the Program shall
              continue and survive.
              Everyone is permitted to copy and distribute copies of this Agreement, but in order to
              avoid inconsistency the Agreement is copyrighted and may only be modified in the
              following manner. The Agreement Steward reserves the right to publish new versions
APPENDIX F — Third-Party Software Licenses                                     RedHawk User Manual | 964
(including revisions) of this Agreement from time to time. No one other than the
              Agreement Steward has the right to modify this Agreement. IBM is the initial Agreement
              Steward. IBM may assign the responsibility to serve as the Agreement Steward to a
              suitable separate entity. Each new version of the Agreement will be given a distinguishing
              version number. The Program (including Contributions) may always be distributed subject
              to the version of the Agreement under which it was received. In addition, after a new
              version of the Agreement is published, Contributor may elect to distribute the Program
              (including its Contributions) under the new version. Except as expressly stated in
              Sections 2(a) and 2(b) above, Recipient receives no rights or licenses to the intellectual
              property of any Contributor under this Agreement, whether expressly, by implication,
              estoppel or otherwise. All rights in the Program not expressly granted under this
              Agreement are reserved.
              This Agreement is governed by the laws of the State of New York and the intellectual
              property laws of the United States of America. No party to this Agreement will bring a
              legal action under this Agreement more than one year after the cause of action arose.
              Each party waives its rights to a jury trial in any resulting litigation.
\end{lstlisting}

\subsection{The MIT License (MIT)}

\begin{lstlisting}[basicstyle=\ttfamily\scriptsize,breaklines=true]
The MIT License (MIT)
              Copyright (c) 2017 ZHU Kaidi
              Permission is hereby granted, free of charge, to any person obtaining a copy of this
              software and associated documentation files (the "Software"), to deal in the Software
              without restriction, including without limitation the rights to use, copy, modify, merge,
              publish, distribute, sublicense, and/or sell copies of the Software, and to permit persons to
              whom the Software is furnished to do so, subject to the following conditions:
              The above copyright notice and this permission notice shall be included in all copies or
              substantial portions of the Software.
              THE SOFTWARE IS PROVIDED "AS IS", WITHOUT WARRANTY OF ANY KIND,
              EXPRESS OR IMPLIED, INCLUDING BUT NOT LIMITED TO THE WARRANTIES OF
              MERCHANTABILITY, FITNESS FOR A PARTICULAR PURPOSE AND
              NONINFRINGEMENT. IN NO EVENT SHALL THE AUTHORS OR COPYRIGHT
              HOLDERS BE LIABLE FOR ANY CLAIM, DAMAGES OR OTHER LIABILITY, WHETHER
              IN AN ACTION OF CONTRACT, TORT OR OTHERWISE, ARISING FROM, OUT OF OR
              IN CONNECTION WITH THE SOFTWARE OR THE USE OR OTHER DEALINGS IN
              THE SOFTWARE.
\end{lstlisting}

\subsection{Permission is hereby granted, free of charge, to any person or organization obtaining a}

\begin{lstlisting}[basicstyle=\ttfamily\scriptsize,breaklines=true]
Permission is hereby granted, free of charge, to any person or organization obtaining a
              copy of the software and accompanying documentation covered by this license (the
              "Software") to use, reproduce, display, distribute, execute, and transmit the Software, and
              to prepare derivative works of the Software, and to permit third-parties to whom the
              Software is furnished to do so, all subject to the following:
              The copyright notices in the Software and this entire statement, including the above
              license grant, this restriction and the following disclaimer, must be included in all copies of
              the Software, in whole or in part, and all derivative works of the Software, unless such
              copies or derivative works are solely in the form of machine-executable object code
              generated by a source language processor.
              THE SOFTWARE IS PROVIDED "AS IS", WITHOUT WARRANTY OF ANY KIND,
              EXPRESS OR IMPLIED, INCLUDING BUT NOT LIMITED TO THE WARRANTIES OF
              MERCHANTABILITY, FITNESS FOR A PARTICULAR PURPOSE, TITLE AND NON-
APPENDIX F — Third-Party Software Licenses                                    RedHawk User Manual | 965
INFRINGEMENT. IN NO EVENT SHALL THE COPYRIGHT HOLDERS OR ANYONE
              DISTRIBUTING THE SOFTWARE BE LIABLE FOR ANY DAMAGES OR OTHER
              LIABILITY, WHETHER IN CONTRACT, TORT OR OTHERWISE, ARISING FROM, OUT
              OF OR IN CONNECTION WITH THE SOFTWARE OR THE USE OR OTHER
              DEALINGS IN THE SOFTWARE.
\end{lstlisting}

\subsection{Copyright (C) 1996-2002 by Jorn Lind-Nielsen All rights reserved}

\begin{lstlisting}[basicstyle=\ttfamily\scriptsize,breaklines=true]
Copyright (C) 1996-2002 by Jorn Lind-Nielsen All rights reserved
              Permission is hereby granted, without written agreement and without license or royalty
              fees, to use, reproduce, prepare derivative works, distribute, and display this software and
              its documentation for any purpose, provided that (1) the above copyright notice and the
              following two paragraphs appear in all copies of the source code and (2) redistributions,
              including without limitation binaries, reproduce these notices in the supporting
              documentation. Substantial modifications to this software may be copyrighted by their
              authors and need not follow the licensing terms described here, provided that the new
              terms are clearly indicated in all files where they apply.
              IN NO EVENT SHALL JORN LIND-NIELSEN, OR DISTRIBUTORS OF THIS
              SOFTWARE BE LIABLE TO ANY PARTY FOR DIRECT, INDIRECT, SPECIAL,
              INCIDENTAL, OR CONSEQUENTIAL DAMAGES ARISING OUT OF THE USE OF THIS
              SOFTWARE AND ITS DOCUMENTATION, EVEN IF THE AUTHORS OR ANY OF THE
              ABOVE PARTIES HAVE BEEN ADVISED OF THE POSSIBILITY OF SUCH DAMAGE.
              JORN LIND-NIELSEN SPECIFICALLY DISCLAIM ANY WARRANTIES, INCLUDING,
              BUT NOT LIMITED TO, THE IMPLIED WARRANTIES OF MERCHANTABILITY AND
              FITNESS FOR A PARTICULAR PURPOSE. THE SOFTWARE PROVIDED
              HEREUNDER IS ON AN "AS IS" BASIS, AND THE AUTHORS AND DISTRIBUTORS
              HAVE NO OBLIGATION TO PROVIDE MAINTENANCE, SUPPORT, UPDATES,
              ENHANCEMENTS, OR MODIFICATIONS.
The GAlib source code is not in the public domain, but it is available at no cost. As a work
              developed by Matthew Wall using MIT resources and MIT funding, the original GAlib
              source code copyright is owned by the Massachusetts Institute of Technology. Later
              portions were added by Matthew Wall without MIT funding. All rights are reserved.
              GAlib License Agreement
              Permission is hereby granted, free of charge, to any person obtaining a copy of this
              software and associated documentation files (the "Software"), to deal in the Software
              without restriction, including without limitation the rights to use, copy, modify, merge,
              publish, distribute, sublicense, and/or sell copies of the Software, and to permit persons to
              whom the Software is furnished to do so, subject to the following conditions:
              You may copy and distribute the source code and/or library/executable code for GAlib in
              any medium provided that you conspicuously and appropriately give credit to the author
              and keep intact all copyright and disclaimer notices in the library.
              Any publications of work based upon experiments that use GAlib must include a suitable
              acknowledgement of GAlib. A suggested acknowledgement is: "The software for this work
              used the GAlib genetic algorithm package, written by Matthew Wall at the Massachusetts
              Institute of Technology."
              The author of GAlib and MIT assume absolutely no responsibility for the use or misuse of
              GAlib. In no event shall the author of GAlib or MIT be liable for any damages resulting
              from use or performance of GAlib.
ANSYS, Inc.

APPENDIX F — Third-Party Software Licenses                                     RedHawk User Manual | 966
\end{lstlisting}

\subsection{The GIFLIB distribution is Copyright (c) 1997 Eric S. Raymond}

\begin{lstlisting}[basicstyle=\ttfamily\scriptsize,breaklines=true]
The GIFLIB distribution is Copyright (c) 1997 Eric S. Raymond
              Permission is hereby granted, free of charge, to any person obtaining a copy of this
              software and associated documentation files (the "Software"), to deal in the Software
              without restriction, including without limitation the rights to use, copy, modify, merge,
              publish, distribute, sublicense, and/or sell copies of the Software, and to permit persons to
              whom the Software is furnished to do so, subject to the following conditions:
              The above copyright notice and this permission notice shall be included in all copies or
              substantial portions of the Software.
              THE SOFTWARE IS PROVIDED "AS IS", WITHOUT WARRANTY OF ANY KIND,
              EXPRESS OR IMPLIED, INCLUDING BUT NOT LIMITED TO THE WARRANTIES OF
              MERCHANTABILITY, FITNESS FOR A PARTICULAR PURPOSE AND
              NONINFRINGEMENT. IN NO EVENT SHALL THE AUTHORS OR COPYRIGHT
              HOLDERS BE LIABLE FOR ANY CLAIM, DAMAGES OR OTHER LIABILITY, WHETHER
              IN AN ACTION OF CONTRACT, TORT OR OTHERWISE, ARISING FROM, OUT OF OR
              IN CONNECTION WITH THE SOFTWARE OR THE USE OR OTHER DEALINGS IN
              THE SOFTWARE.
\end{lstlisting}

\subsection{Copyright (c) 1992-2013 The University of Tennessee and The University of Tennessee}

\begin{lstlisting}[basicstyle=\ttfamily\scriptsize,breaklines=true]
Copyright (c) 1992-2013 The University of Tennessee and The University of Tennessee
              Research Foundation. All rights reserved. Copyright (c) 2000-2013 The University of
              California Berkeley. All rights reserved. Copyright (c) 2006-2013 The University of
              Colorado Denver. All rights reserved.
              $COPYRIGHT$
              Additional copyrights may follow
              $HEADER$
              Redistribution and use in source and binary forms, with or without modification, are
              permitted provided that the following conditions are met:
              - Redistributions of source code must retain the above copyright notice, this list of
              conditions and the following disclaimer.
              - Redistributions in binary form must reproduce the above copyright notice, this list of
              conditions and the following disclaimer listed in this license in the documentation and/or
              other materials provided with the distribution.
              - Neither the name of the copyright holders nor the names of its contributors may be
              used to endorse or promote products derived from this software without specific prior
              written permission.
              The copyright holders provide no reassurances that the source code provided does not
              infringe any patent, copyright, or any other intellectual property rights of third parties. The
              copyright holders disclaim any liability to any recipient for claims brought against recipient
              by any third party for infringement of that parties intellectual property rights.
              THIS SOFTWARE IS PROVIDED BY THE COPYRIGHT HOLDERS AND
              CONTRIBUTORS "AS IS" AND ANY EXPRESS OR IMPLIED WARRANTIES,
              INCLUDING, BUT NOT LIMITED TO, THE IMPLIED WARRANTIES OF
              MERCHANTABILITY AND FITNESS FOR A PARTICULAR PURPOSE ARE
              DISCLAIMED. IN NO EVENT SHALL THE COPYRIGHT OWNER OR CONTRIBUTORS
              BE LIABLE FOR ANY DIRECT, INDIRECT, INCIDENTAL, SPECIAL, EXEMPLARY, OR
              CONSEQUENTIAL DAMAGES (INCLUDING, BUT NOT LIMITED TO, PROCUREMENT
              OF SUBSTITUTE GOODS OR SERVICES; LOSS OF USE, DATA, OR PROFITS; OR
APPENDIX F — Third-Party Software Licenses                                   RedHawk User Manual | 967
BUSINESS INTERRUPTION) HOWEVER CAUSED AND ON ANY THEORY OF
              LIABILITY, WHETHER IN CONTRACT, STRICT LIABILITY, OR TORT (INCLUDING
              NEGLIGENCE OR OTHERWISE) ARISING IN ANY WAY OUT OF THE USE OF THIS
              SOFTWARE, EVEN IF ADVISED OF THE POSSIBILITY OF SUCH DAMAGE.
\end{lstlisting}

\subsection{Copyright (c) 2003-2009 Tim Kientzle All rights reserved.}

\begin{lstlisting}[basicstyle=\ttfamily\scriptsize,breaklines=true]
Copyright (c) 2003-2009 Tim Kientzle All rights reserved.
              Redistribution and use in source and binary forms, with or without modification, are
              permitted provided that the following conditions are met: 1. Redistributions of source code
              must retain the above copyright notice, this list of conditions and the following disclaimer
              in this position and unchanged. 2. Redistributions in binary form must reproduce the
              above copyright notice, this list of conditions and the following disclaimer in the
              documentation and/or other materials provided with the distribution.
              THIS SOFTWARE IS PROVIDED BY THE AUTHOR(S) ``AS IS'' AND ANY EXPRESS
              OR IMPLIED WARRANTIES, INCLUDING, BUT NOT LIMITED TO, THE IMPLIED
              WARRANTIES OF MERCHANTABILITY AND FITNESS FOR A PARTICULAR
              PURPOSE ARE DISCLAIMED. IN NO EVENT SHALL THE AUTHOR(S) BE LIABLE
              FOR ANY DIRECT, INDIRECT, INCIDENTAL, SPECIAL, EXEMPLARY, OR
              CONSEQUENTIAL DAMAGES (INCLUDING, BUT NOT LIMITED TO, PROCUREMENT
              OF SUBSTITUTE GOODS OR SERVICES; LOSS OF USE, DATA, OR PROFITS; OR
              BUSINESS INTERRUPTION) HOWEVER CAUSED AND ON ANY THEORY OF
              LIABILITY, WHETHER IN CONTRACT, STRICT LIABILITY, OR TORT (INCLUDING
              NEGLIGENCE OR OTHERWISE) ARISING IN ANY WAY OUT OF THE USE OF THIS
              SOFTWARE, EVEN IF ADVISED OF THE POSSIBILITY OF SUCH DAMAGE.
\end{lstlisting}

\subsection{Copyright (C) 2002-2011 Carsten Haitzler and various contributors (see AUTHORS)}

\begin{lstlisting}[basicstyle=\ttfamily\scriptsize,breaklines=true]
Copyright (C) 2002-2011 Carsten Haitzler and various contributors (see AUTHORS)
              All rights reserved.
              Redistribution and use in source and binary forms, with or without modification, are
              permitted provided that the following conditions are met:
                1. Redistributions of source code must retain the above copyright notice, this list of
              conditions and the following disclaimer.
                2. Redistributions in binary form must reproduce the above copyright notice, this list of
              conditions and the following disclaimer in the documentation and/or other materials
              provided with the distribution.
              THIS SOFTWARE IS PROVIDED "AS IS" AND ANY EXPRESS OR IMPLIED
              WARRANTIES, INCLUDING, BUT NOT LIMITED TO, THE IMPLIED WARRANTIES OF
              MERCHANTABILITY AND FITNESS FOR A PARTICULAR PURPOSE ARE
              DISCLAIMED. IN NO EVENT SHALL THE COPYRIGHT HOLDER OR CONTRIBUTORS
              BE LIABLE FOR ANY DIRECT, INDIRECT, INCIDENTAL, SPECIAL, EXEMPLARY, OR
              CONSEQUENTIAL DAMAGES (INCLUDING, BUT NOT LIMITED TO, PROCUREMENT
              OF SUBSTITUTE GOODS OR SERVICES; LOSS OF USE, DATA, OR PROFITS; OR
              BUSINESS INTERRUPTION) HOWEVER CAUSED AND ON ANY THEORY OF
              LIABILITY, WHETHER IN CONTRACT, STRICT LIABILITY, OR TORT (INCLUDING
              NEGLIGENCE OR OTHERWISE) ARISING IN ANY WAY OUT OF THE USE OF THIS
              SOFTWARE, EVEN IF ADVISED OF THE POSSIBILITY OF SUCH DAMAGE.
ANSYS, Inc.

APPENDIX F — Third-Party Software Licenses                                     RedHawk User Manual | 968
Copyright 1990 - 1997 by AT\&T, Lucent Technologies and Bellcore.
              Permission to use, copy, modify, and distribute this software and its documentation for
              any purpose and without fee is hereby granted, provided that the above copyright notice
              appear in all copies and that both that the copyright notice and this permission notice and
              warranty disclaimer appear in supporting documentation, and that the names of AT\&T,
              Bell Laboratories, Lucent or Bellcore or any of their entities not be used in advertising or
              publicity pertaining to distribution of the software without specific, written prior permission.
              AT\&T, Lucent and Bellcore disclaim all warranties with regard to this software, including
              all implied warranties of merchantability and fitness. In no event shall AT\&T, Lucent or
              Bellcore be liable for any special, indirect or consequential damages or any damages
              whatsoever resulting from loss of use, data or profits, whether in an action of contract,
              negligence or other tortious action, arising out of or in connection with the use or
              performance of this software.
\end{lstlisting}

\subsection{LIbigl is licensed under the Mozilla Public License (MPL) Version 2.0, the text of which}

\begin{lstlisting}[basicstyle=\ttfamily\scriptsize,breaklines=true]
LIbigl is licensed under the Mozilla Public License (MPL) Version 2.0, the text of which
              can be found at: https://www.mozilla.org/media/MPL/2.0/index.815ca599c9df.txt. Please
              contact open.source@ansys.com for a copy of the Libigl source code.
This software is based in part on the work of the Independent JPEG Group.
\end{lstlisting}

\subsection{Licensed under the Apache License, Version 2.0 (the "License"); you may not use this file}

\begin{lstlisting}[basicstyle=\ttfamily\scriptsize,breaklines=true]
Licensed under the Apache License, Version 2.0 (the "License"); you may not use this file
              except in compliance with the License.
You may obtain a copy of the License at
              http://www.apache.org/licenses/LICENSE-2.0
Unless required by applicable law or agreed to in writing, software distributed under the
              License is distributed on an "AS IS" BASIS, WITHOUT WARRANTIES OR CONDITIONS
              OF ANY KIND, either express or implied. See the License for the specific language
              governing permissions and limitations under the License.
Redistribution and use in source and binary forms, with or without modification, are
              permitted provided that the following conditions are met:
1. Redistributions of source code must retain the above copyright notice, this list of
              conditions and the following disclaimer.
ANSYS, Inc.

APPENDIX F — Third-Party Software Licenses                                   RedHawk User Manual | 969
2. Redistributions in binary form must reproduce the above copyright notice, this list of
              conditions and the following disclaimer in the documentation and/or other materials
              provided with the distribution.
3. Neither the name of the project nor the names of its contributors may be used to
              endorse or promote products derived from this software without specific prior written
              permission.
THIS SOFTWARE IS PROVIDED BY THE COPYRIGHT HOLDERS AND
              CONTRIBUTORS "AS IS" AND ANY EXPRESS OR IMPLIED WARRANTIES,
              INCLUDING, BUT NOT LIMITED TO, THE IMPLIED WARRANTIES OF
              MERCHANTABILITY AND FITNESS FOR A PARTICULAR PURPOSE ARE
              DISCLAIMED. IN NO EVENT SHALL THE COPYRIGHT OWNER OR CONTRIBUTORS
              BE LIABLE FOR ANY DIRECT, INDIRECT, INCIDENTAL, SPECIAL, EXEMPLARY, OR
              CONSEQUENTIAL DAMAGES (INCLUDING, BUT NOT LIMITED TO, PROCUREMENT
              OF SUBSTITUTE GOODS OR SERVICES; LOSS OF USE, DATA, OR PROFITS; OR
              BUSINESS INTERRUPTION) HOWEVER CAUSED AND ON ANY THEORY OF
              LIABILITY, WHETHER IN CONTRACT, STRICT LIABILITY, OR TORT (INCLUDING
              NEGLIGENCE OR OTHERWISE) ARISING IN ANY WAY OUT OF THE USE OF THIS
              SOFTWARE, EVEN IF ADVISED OF THE POSSIBILITY OF SUCH DAMAGE.
\end{lstlisting}

\subsection{Copyright (c) 2005 Nominet UK (www.nic.uk)}

\begin{lstlisting}[basicstyle=\ttfamily\scriptsize,breaklines=true]
Copyright (c) 2005 Nominet UK (www.nic.uk)
              All rights reserved.
Contributors: Ben Laurie, Rachel Willmer. The Contributors have asserted their moral
              rights under the UK Copyright Design and Patents Act 1988 to be recorded as the authors
              of this copyright work.
              This is a BSD-style Open Source licence.
Redistribution and use in source and binary forms, with or without modification, are
              permitted provided that the following conditions are met:
1. Redistributions of source code must retain the above copyright notice, this list of
              conditions and the following disclaimer.
2. Redistributions in binary form must reproduce the above copyright notice, this list of
              conditions and the following disclaimer in the documentation and/or other materials
              provided with the distribution.
3. The name of Nominet UK or the contributors may not be used to endorse or promote
              products derived from this software without specific prior written permission;
              and provided that the user accepts the terms of the following disclaimer:
THIS SOFTWARE IS PROVIDED BY NOMINET UK AND CONTRIBUTORS ``AS IS''
              AND ANY EXPRESS OR IMPLIED WARRANTIES, INCLUDING, BUT NOT LIMITED TO,
APPENDIX F — Third-Party Software Licenses                                   RedHawk User Manual | 970
THE IMPLIED WARRANTIES OF MERCHANTABILITY AND FITNESS FOR A
              PARTICULAR PURPOSE ARE DISCLAIMED. IN NO EVENT SHALL NOMINET UK OR
              CONTRIBUTORS BE LIABLE FOR ANY DIRECT, INDIRECT, INCIDENTAL, SPECIAL,
              EXEMPLARY, OR CONSEQUENTIAL DAMAGES (INCLUDING, BUT NOT LIMITED TO,
              PROCUREMENT OF SUBSTITUTE GOODS OR SERVICES; LOSS OF USE, DATA, OR
              PROFITS; OR BUSINESS INTERRUPTION) HOWEVER CAUSED AND ON ANY
              THEORY OF LIABILITY, WHETHER IN CONTRACT, STRICT LIABILITY, OR TORT
              (INCLUDING NEGLIGENCE OR OTHERWISE) ARISING IN ANY WAY OUT OF THE
              USE OF THIS SOFTWARE, EVEN IF ADVISED OF THE POSSIBILITY OF SUCH
              DAMAGE.
\end{lstlisting}

\subsection{Copyright (c) 1998-2008 The OpenSSL Project. All rights reserved.}

\begin{lstlisting}[basicstyle=\ttfamily\scriptsize,breaklines=true]
Copyright (c) 1998-2008 The OpenSSL Project. All rights reserved.
Redistribution and use in source and binary forms, with or without modification, are
              permitted provided that the following conditions are met:
1. Redistributions of source code must retain the above copyright notice, this list of
              conditions and the following disclaimer.
2. Redistributions in binary form must reproduce the above copyright notice, this list of
              conditions and the following disclaimer in the documentation and/or other materials
              provided with the distribution.
3. All advertising materials mentioning features or use of this software must display the
              following acknowledgment: "This product includes software developed by the OpenSSL
              Project for use in the OpenSSL Toolkit. (http://www.openssl.org/)"
4. The names "OpenSSL Toolkit" and "OpenSSL Project" must not be used to endorse or
              promote products derived from this software without prior written permission. For written
              permission, please contact openssl-core@openssl.org.
5. Products derived from this software may not be called "OpenSSL" nor may "OpenSSL"
              appear in their names without prior written permission of the OpenSSL Project.
6. Redistributions of any form whatsoever must retain the following acknowledgment:
              "This product includes software developed by the OpenSSL Project for use in the
              OpenSSL Toolkit (http://www.openssl.org/)"
THIS SOFTWARE IS PROVIDED BY THE OpenSSL PROJECT ``AS IS'' AND ANY
              EXPRESSED OR IMPLIED WARRANTIES, INCLUDING, BUT NOT LIMITED TO, THE
              IMPLIED WARRANTIES OF MERCHANTABILITY AND FITNESS FOR A PARTICULAR
              PURPOSE ARE DISCLAIMED. IN NO EVENT SHALL THE OpenSSL PROJECT OR ITS
              CONTRIBUTORS BE LIABLE FOR ANY DIRECT, INDIRECT, INCIDENTAL, SPECIAL,
              EXEMPLARY, OR CONSEQUENTIAL DAMAGES (INCLUDING, BUT NOT LIMITED TO,
              PROCUREMENT OF SUBSTITUTE GOODS OR SERVICES; LOSS OF USE, DATA, OR
              PROFITS; OR BUSINESS INTERRUPTION) HOWEVER CAUSED AND ON ANY
              THEORY OF LIABILITY, WHETHER IN CONTRACT, STRICT LIABILITY, OR TORT
APPENDIX F — Third-Party Software Licenses                                   RedHawk User Manual | 971
(INCLUDING NEGLIGENCE OR OTHERWISE) ARISING IN ANY WAY OUT OF THE
              USE OF THIS SOFTWARE, EVEN IF ADVISED OF THE POSSIBILITY OF SUCH
              DAMAGE.
              ===================================================================
This product includes cryptographic software written by Eric Young (eay@cryptsoft.com).
              This product includes software written by Tim Hudson (tjh@cryptsoft.com).
\end{lstlisting}

\subsection{Copyright (C) 1995-1998 Eric Young (eay@cryptsoft.com) All rights reserved.}

\begin{lstlisting}[basicstyle=\ttfamily\scriptsize,breaklines=true]
Copyright (C) 1995-1998 Eric Young (eay@cryptsoft.com) All rights reserved.
This package is an SSL implementation written by Eric Young (eay@cryptsoft.com). The
              implementation was written so as to conform with Netscapes SSL.
This library is free for commercial and non-commercial use as long as the following
              conditions are aheared to. The following conditions apply to all code found in this
              distribution, be it the RC4, RSA, lhash, DES, etc., code; not just the SSL code. The SSL
              documentation included with this distribution is covered by the same copyright terms
              except that the holder is Tim Hudson (tjh@cryptsoft.com).
Copyright remains Eric Young's, and as such any Copyright notices in the code are not to
              be removed. If this package is used in a product, Eric Young should be given attribution
              as the author of the parts of the library used. This can be in the form of a textual message
              at program startup or in documentation (online or textual) provided with the package.
Redistribution and use in source and binary forms, with or without modification, are
              permitted provided that the following conditions are met: 1. Redistributions of source code
              must retain the copyright notice, this list of conditions and the following disclaimer. 2.
              Redistributions in binary form must reproduce the above copyright notice, this list of
              conditions and the following disclaimer in the documentation and/or other materials
              provided with the distribution. 3. All advertising materials mentioning features or use of
              this software must display the following acknowledgement: "This product includes
              cryptographic software written by Eric Young (eay@cryptsoft.com)" The word
              'cryptographic' can be left out if the rouines from the library being used are not
              cryptographic related :-). 4. If you include any Windows specific code (or a derivative
              thereof) from the apps directory (application code) you must include an
              acknowledgement: "This product includes software written by Tim Hudson
              (tjh@cryptsoft.com)"
THIS SOFTWARE IS PROVIDED BY ERIC YOUNG ``AS IS'' AND ANY EXPRESS OR
              IMPLIED WARRANTIES, INCLUDING, BUT NOT LIMITED TO, THE IMPLIED
              WARRANTIES OF MERCHANTABILITY AND FITNESS FOR A PARTICULAR
              PURPOSE ARE DISCLAIMED. IN NO EVENT SHALL THE AUTHOR OR
              CONTRIBUTORS BE LIABLE FOR ANY DIRECT, INDIRECT, INCIDENTAL, SPECIAL,
              EXEMPLARY, OR CONSEQUENTIAL DAMAGES (INCLUDING, BUT NOT LIMITED TO,
              PROCUREMENT OF SUBSTITUTE GOODS OR SERVICES; LOSS OF USE, DATA, OR
              PROFITS; OR BUSINESS INTERRUPTION) HOWEVER CAUSED AND ON ANY
              THEORY OF LIABILITY, WHETHER IN CONTRACT, STRICT LIABILITY, OR TORT
              (INCLUDING NEGLIGENCE OR OTHERWISE) ARISING IN ANY WAY OUT OF THE
APPENDIX F — Third-Party Software Licenses                                    RedHawk User Manual | 972
USE OF THIS SOFTWARE, EVEN IF ADVISED OF THE POSSIBILITY OF SUCH
              DAMAGE.
The licence and distribution terms for any publically available version or derivative of this
              code cannot be changed. i.e. this code cannot simply be copied and put under another
              distribution licence [including the GNU Public Licence.]
Redistribution and use in source and binary forms, with or without
              modification, are permitted provided that the following conditions are met:
Redistributions of source code must retain the above copyright notice, this list of
              conditions and the following disclaimer.
Redistributions in binary form must reproduce the above copyright notice, this list of
              conditions and the following disclaimer in the documentation and/or other materials
              provided with the distribution.
Neither the name of the University of Cambridge nor the name of Google Inc. nor the
              names of their contributors may be used to endorse or promote products derived from this
              software without specific prior written permission.
THIS SOFTWARE IS PROVIDED BY THE COPYRIGHT HOLDERS AND
              CONTRIBUTORS "AS IS"AND ANY EXPRESS OR IMPLIED WARRANTIES,
              INCLUDING, BUT NOT LIMITED TO, THEIMPLIED WARRANTIES OF
              MERCHANTABILITY AND FITNESS FOR A PARTICULAR PURPOSEARE
              DISCLAIMED. IN NO EVENT SHALL THE COPYRIGHT OWNER OR CONTRIBUTORS
              BELIABLE FOR ANY DIRECT, INDIRECT, INCIDENTAL, SPECIAL, EXEMPLARY,
              ORCONSEQUENTIAL DAMAGES (INCLUDING, BUT NOT LIMITED TO,
              PROCUREMENT OFSUBSTITUTE GOODS OR SERVICES; LOSS OF USE, DATA, OR
              PROFITS; OR BUSINESSINTERRUPTION) HOWEVER CAUSED AND ON ANY
              THEORY OF LIABILITY, WHETHER INCONTRACT, STRICT LIABILITY, OR TORT
              (INCLUDING NEGLIGENCE OR OTHERWISE)ARISING IN ANY WAY OUT OF THE
              USE OF THIS SOFTWARE, EVEN IF ADVISED OF THEPOSSIBILITY OF SUCH
              DAMAGE.
\end{lstlisting}

\subsection{Copyright (c) 2009, Krasnoshchekov Petr}

\begin{lstlisting}[basicstyle=\ttfamily\scriptsize,breaklines=true]
Copyright (c) 2009, Krasnoshchekov Petr
              All rights reserved.
Redistribution and use in source and binary forms, with or without modification, are
              permitted provided that the following conditions are met:
Redistributions of source code must retain the above copyright notice, this list of
              conditions and the following disclaimer.
ANSYS, Inc.

APPENDIX F — Third-Party Software Licenses                                   RedHawk User Manual | 973
Redistributions in binary form must reproduce the above copyright notice, this list of
              conditions and the following disclaimer in the documentation and/or other materials
              provided with the distribution.
Neither the name of the <organization> nor the names of its contributors may be used to
              endorse or promote products derived from this software without specific prior written
              permission.
THIS SOFTWARE IS PROVIDED BY Krasnoshchekov Petr ''AS IS'' AND ANY EXPRESS
              OR IMPLIED WARRANTIES, INCLUDING, BUT NOT LIMITED TO, THE IMPLIED
              WARRANTIES OF MERCHANTABILITY AND FITNESS FOR A PARTICULAR
              PURPOSE ARE DISCLAIMED. IN NO EVENT SHALL Krasnoshchekov Petr BE LIABLE
              FOR ANY DIRECT, INDIRECT, INCIDENTAL, SPECIAL, EXEMPLARY, OR
              CONSEQUENTIAL DAMAGES (INCLUDING, BUT NOT LIMITED TO, PROCUREMENT
              OF SUBSTITUTE GOODS OR SERVICES; LOSS OF USE, DATA, OR PROFITS; OR
              BUSINESS INTERRUPTION) HOWEVER CAUSED AND ON ANY THEORY OF
              LIABILITY, WHETHER IN CONTRACT, STRICT LIABILITY, OR TORT (INCLUDING
              NEGLIGENCE OR OTHERWISE) ARISING IN ANY WAY OUT OF THE USE OF THIS
              SOFTWARE, EVEN IF ADVISED OF THE POSSIBILITY OF SUCH DAMAGE.
\end{lstlisting}

\subsection{Copyright [Copyright (c) 1994-1996 The Regents of the Univ. of California. All rights}

\begin{lstlisting}[basicstyle=\ttfamily\scriptsize,breaklines=true]
Copyright [Copyright (c) 1994-1996 The Regents of the Univ. of California. All rights
              reserved.
\end{lstlisting}

\subsection{Permission is hereby granted, without written agreement and without license or royalty}

\begin{lstlisting}[basicstyle=\ttfamily\scriptsize,breaklines=true]
Permission is hereby granted, without written agreement and without license or royalty
              fees, to use, copy, modify, and distribute this software and its documentation for any
              purpose, provided that the above copyright notice and the following two paragraphs
              appear in all copies of this software.
IN NO EVENT SHALL THE UNIVERSITY OF CALIFORNIA BE LIABLE TO ANY PARTY
              FOR DIRECT, INDIRECT, SPECIAL, INCIDENTAL, OR CONSEQUENTIAL DAMAGES
              ARISING OUT OF THE USE OF THIS SOFTWARE AND ITS DOCUMENTATION, EVEN
              IF THE UNIVERSITY OF CALIFORNIA HAS BEEN ADVISED OF THE POSSIBILITY OF
              SUCH DAMAGE.
THE UNIVERSITY OF CALIFORNIA SPECIFICALLY DISCLAIMS ANY WARRANTIES,
              INCLUDING, BUT NOT LIMITED TO, THE IMPLIED WARRANTIES OF
              MERCHANTABILITY AND FITNESS FOR A PARTICULAR PURPOSE. THE
              SOFTWARE PROVIDED HEREUNDER IS ON AN "AS IS" BASIS, AND THE
              UNIVERSITY OF CALIFORNIA HAS NO OBLIGATION TO PROVIDE MAINTENANCE,
              SUPPORT, UPDATES, ENHANCEMENTS, OR MODIFICATIONS.]
This program is based in part on the work of the Qwt project (http://qwt.sf.net).
\end{lstlisting}

\subsection{Copyright (c) 1985,86,87,88 by Kenneth S. Kundert and the University of California.}

\begin{lstlisting}[basicstyle=\ttfamily\scriptsize,breaklines=true]
Copyright (c) 1985,86,87,88 by Kenneth S. Kundert and the University of California.
APPENDIX F — Third-Party Software Licenses                                     RedHawk User Manual | 974
Permission to use, copy, modify, and distribute this software and its documentation for
              any purpose and without fee is hereby granted, provided that the copyright notices appear
              in all copies and supporting documentation and that the authors and the University of
              California are properly credited. The authors and the University of California make no
              representations as to the suitability of this software for any purpose. It is provided `as is',
              without express or implied warranty.
\end{lstlisting}

\subsection{Copyright (c) 2003, The Regents of the University of California, through Lawrence}

\begin{lstlisting}[basicstyle=\ttfamily\scriptsize,breaklines=true]
Copyright (c) 2003, The Regents of the University of California, through Lawrence
              Berkeley National Laboratory (subject to receipt of any required approvals from U.S.
              Dept. of Energy)
Redistribution and use in source and binary forms, with or without modification, are
              permitted provided that the following conditions are met:
(1) Redistributions of source code must retain the above copyright notice, this list of
              conditions and the following disclaimer. (2) Redistributions in binary form must reproduce
              the above copyright notice, this list of conditions and the following disclaimer in the
              documentation and/or other materials provided with the distribution. (3) Neither the name
              of Lawrence Berkeley National Laboratory, U.S. Dept. of Energy nor the names of its
              contributors may be used to endorse or promote products derived from this software
              without specific prior written permission.
THIS SOFTWARE IS PROVIDED BY THE COPYRIGHT HOLDERS AND
              CONTRIBUTORS "AS IS" AND ANY EXPRESS OR IMPLIED WARRANTIES,
              INCLUDING, BUT NOT LIMITED TO, THE IMPLIED WARRANTIES OF
              MERCHANTABILITY AND FITNESS FOR A PARTICULAR PURPOSE ARE
              DISCLAIMED. IN NO EVENT SHALL THE COPYRIGHT OWNER OR CONTRIBUTORS
              BE LIABLE FOR ANY DIRECT, INDIRECT, INCIDENTAL, SPECIAL, EXEMPLARY, OR
              CONSEQUENTIAL DAMAGES (INCLUDING, BUT NOT LIMITED TO, PROCUREMENT
              OF SUBSTITUTE GOODS OR SERVICES; LOSS OF USE, DATA, OR PROFITS; OR
              BUSINESS INTERRUPTION) HOWEVER CAUSED AND ON ANY THEORY OF
              LIABILITY, WHETHER IN CONTRACT, STRICT LIABILITY, OR TORT (INCLUDING
              NEGLIGENCE OR OTHERWISE) ARISING IN ANY WAY OUT OF THE USE OF THIS
              SOFTWARE, EVEN IF ADVISED OF THE POSSIBILITY OF SUCH DAMAGE.
Redistribution and use in source and binary forms, with or without modification, are
              permitted provided that the following conditions are met:
ANSYS, Inc.

APPENDIX F — Third-Party Software Licenses                                   RedHawk User Manual | 975
1. Redistributions of source code must retain the above copyright notice, this list of
              conditions and the following disclaimer.
2. Redistributions in binary form must reproduce the above copyright notice, this list of
              conditions and the following disclaimer in the documentation and/or other materials
              provided with the distribution.
3. The names of its contributors may not be used to endorse or promote      products
              derived from this software without specific prior written permission.
THIS SOFTWARE IS PROVIDED BY THE COPYRIGHT HOLDERS AND
              CONTRIBUTORS "AS IS" AND ANY EXPRESS OR IMPLIED WARRANTIES,
              INCLUDING, BUT NOT LIMITED TO, THE IMPLIED WARRANTIES OF
              MERCHANTABILITY AND FITNESS FOR A PARTICULAR PURPOSE ARE
              DISCLAIMED. IN NO EVENT SHALL THE COPYRIGHT OWNER OR CONTRIBUTORS
              BE LIABLE FOR ANY DIRECT, INDIRECT, INCIDENTAL, SPECIAL, EXEMPLARY, OR
              CONSEQUENTIAL DAMAGES (INCLUDING, BUT NOT LIMITED TO, PROCUREMENT
              OF SUBSTITUTE GOODS OR SERVICES; LOSS OF USE, DATA, OR PROFITS; OR
              BUSINESS INTERRUPTION) HOWEVER CAUSED AND ON ANY THEORY OF
              LIABILITY, WHETHER IN CONTRACT, STRICT LIABILITY, OR TORT (INCLUDING
              NEGLIGENCE OR OTHERWISE) ARISING IN ANY WAY OUT OF THE USE OF THIS
              SOFTWARE, EVEN IF ADVISED OF THE POSSIBILITY OF SUCH DAMAGE.
\end{lstlisting}

\subsection{Copyright, and License:}

\begin{lstlisting}[basicstyle=\ttfamily\scriptsize,breaklines=true]
Copyright, and License:
              Copyright (c) 2005 by Timothy A. Davis, University of Florida, davis@cise.ufl.edu. All
              Rights Reserved.
              UMFPACK License:
              Your use or distribution of UMFPACK or any derivative code implies that you agree to this
              License.
              THIS MATERIAL IS PROVIDED AS IS, WITH ABSOLUTELY NO WARRANTY
              EXPRESSED OR IMPLIED. ANY USE IS AT YOUR OWN RISK.
              Permission is hereby granted to use or copy this program, provided that the Copyright,
              this License, and the Availability of the original version is retained on all copies. User
              documentation of any code that uses UMFPACK or any modified version of UMFPACK
              code must cite the Copyright, this License, the Availability note, and "Used by
              permission." Permission to modify the code and to distribute modified code is granted,
              provided the Copyright, this License, and the Availability note are retained, and a notice
              that the code was modified is included. This software was developed with support from
              the National Science Foundation, and is provided to you free of charge.
Most files in this release are marked with the copyrights of the organizations who have
              edited them. The copyrights below are in no particular order and generally reflect
              members of the Open MPI core team who have contributed code to this release. The
              copyrights for code used under license from other parties are included in the
              corresponding files.
ANSYS, Inc.

APPENDIX F — Third-Party Software Licenses                                    RedHawk User Manual | 976
\end{lstlisting}

\subsection{Copyright (c) 2004-2010 The Trustees of Indiana University and Indiana University}

\begin{lstlisting}[basicstyle=\ttfamily\scriptsize,breaklines=true]
Copyright (c) 2004-2010 The Trustees of Indiana University and Indiana University
              Research and Technology Corporation. All rights reserved.
              Copyright (c) 2004-2010 The University of Tennessee and The University of Tennessee
              Research Foundation. All rights reserved.
              Copyright (c) 2004-2010 High Performance Computing Center Stuttgart, University of
              Stuttgart. All rights reserved.
              Copyright (c) 2004-2008 The Regents of the University of California. All rights reserved.
              Copyright (c) 2006-2010 Los Alamos National Security, LLC. All rights reserved.
              Copyright (c) 2006-2010 Cisco Systems, Inc. All rights reserved.
              Copyright (c) 2006-2010 Voltaire, Inc. All rights reserved.
              Copyright (c) 2006-2011 Sandia National Laboratories. All rights reserved.
              Copyright (c) 2006-2010 Sun Microsystems, Inc. All rights reserved. Use is subject to
              license terms.
              Copyright (c) 2006-2010 The University of Houston. All rights reserved.
              Copyright (c) 2006-2009 Myricom, Inc. All rights reserved.
              Copyright (c) 2007-2008 UT-Battelle, LLC. All rights reserved.
              Copyright (c) 2007-2010 IBM Corporation. All rights reserved.
              Copyright (c) 1998-2005 Forschungszentrum Juelich, Juelich Supercomputing Centre,
              Federal Republic of Germany
              Copyright (c) 2005-2008 ZIH, TU Dresden, Federal Republic of Germany
              Copyright (c) 2007 Evergrid, Inc. All rights reserved.
              Copyright (c) 2008 Chelsio, Inc. All rights reserved.
              Copyright (c) 2008-2009 Institut National de Recherche en Informatique. All rights
              reserved.
              Copyright (c) 2007 Lawrence Livermore National Security, LLC. All rights reserved.
              Copyright (c) 2007-2009 Mellanox Technologies. All rights reserved.
              Copyright (c) 2006-2010 QLogic Corporation. All rights reserved.
              Copyright (c) 2008-2010 Oak Ridge National Labs. All rights reserved.
              Copyright (c) 2006-2010 Oracle and/or its affiliates. All rights reserved.
              Copyright (c) 2009 Bull SAS. All rights reserved.
              Copyright (c) 2010 ARM ltd. All rights reserved.
              Copyright (c) 2010-2011 Alex Brick <bricka (a] ccs.neu.edu>. All rights reserved.
              Copyright (c) 2012 The University of Wisconsin-La Crosse. All rights reserved.
$HEADER$
              Redistribution and use in source and binary forms, with or without modification, are
              permitted provided that the following conditions are met:
              Redistributions of source code must retain the above copyright notice, this list of
              conditions and the following disclaimer.
ANSYS, Inc.

APPENDIX F — Third-Party Software Licenses                                     RedHawk User Manual | 977
Redistributions in binary form must reproduce the above copyright notice, this list of
              conditions and the following disclaimer listed in this license in the documentation and/or
              other materials provided with the distribution.
              Neither the name of the copyright holders nor the names of its contributors may be used
              to endorse or promote products derived from this software without specific prior written
              permission.
              The copyright holders provide no reassurances that the source code provided does not
              infringe any patent, copyright, or any other intellectual property rights of third parties. The
              copyright holders disclaim any liability to any recipient for claims brought against recipient
              by any third party for infringement of that parties intellectual property rights.
THIS SOFTWARE IS PROVIDED BY THE COPYRIGHT HOLDERS AND
              CONTRIBUTORS “AS IS” AND ANY EXPRESS OR IMPLIED WARRANTIES,
              INCLUDING, BUT NOT LIMITED TO, THE IMPLIED WARRANTIES OF
              MERCHANTABILITY AND FITNESS FOR A PARTICULAR PURPOSE ARE
              DISCLAIMED. IN NO EVENT SHALL THE COPYRIGHT OWNER OR CONTRIBUTORS
              BE LIABLE FOR ANY DIRECT, INDIRECT, INCIDENTAL, SPECIAL, EXEMPLARY, OR
              CONSEQUENTIAL DAMAGES (INCLUDING, BUT NOT LIMITED TO, PROCUREMENT
              OF SUBSTITUTE GOODS OR SERVICES; LOSS OF USE, DATA, OR PROFITS; OR
              BUSINESS INTERRUPTION) HOWEVER CAUSED AND ON ANY THEORY OF
              LIABILITY, WHETHER IN CONTRACT, STRICT LIABILITY, OR TORT (INCLUDING
              NEGLIGENCE OR OTHERWISE) ARISING IN ANY WAY OUT OF THE USE OF THIS
              SOFTWARE, EVEN IF ADVISED OF THE POSSIBILITY OF SUCH DAMAGE.
This software is copyrighted by the Regents of the University of California, Sun
              Microsystems, Inc., Scriptics Corporation, and other parties. The following terms apply to
              all files associated with the software, unless explicitly disclaimed in individual files.
The authors hereby grant permission to use, copy, modify, distribute, and license this
              software and its documentation for any purpose, provided that existing copyright notices
              are retained in all copies and that this notice is included verbatim in any distributions. No
              written agreement, license, or royalty fee is required for any of the authorized uses.
              Modifications to this software may be copyrighted by their authors and need not follow the
              licensing terms described here, provided that the new terms are clearly indicated on the
              first page of each file where they apply.
IN NO EVENT SHALL THE AUTHORS OR DISTRIBUTORS BE LIABLE TO ANY PARTY
              FOR DIRECT, INDIRECT, SPECIAL, INCIDENTAL, OR CONSEQUENTIAL DAMAGES
              ARISING OUT OF THE USE OF THIS SOFTWARE, ITS DOCUMENTATION, OR ANY
              DERIVATIVES THEREOF, EVEN IF THE AUTHORS HAVE BEEN ADVISED OF THE
              POSSIBILITY OF SUCH DAMAGE.
              THE AUTHORS AND DISTRIBUTORS SPECIFICALLY DISCLAIM ANY WARRANTIES,
              INCLUDING, BUT NOT LIMITED TO, THE IMPLIED WARRANTIES OF
              MERCHANTABILITY, FITNESS FOR A PARTICULAR PURPOSE, AND NON-
              INFRINGEMENT. THIS SOFTWARE IS PROVIDED ON AN “AS IS” BASIS, AND THE
              AUTHORS AND DISTRIBUTORS HAVE NO OBLIGATION TO PROVIDE
              MAINTENANCE, SUPPORT, UPDATES, ENHANCEMENTS, OR MODIFICATIONS.
ANSYS, Inc.

APPENDIX F — Third-Party Software Licenses                                     RedHawk User Manual | 978
GOVERNMENT USE: If you are acquiring this software on behalf of the U.S. government,
              the Government shall have only “Restricted Rights” in the software and related
              documentation as defined in the Federal Acquisition Regulations (FARs) in Clause
              52.227.19 (c) (2). If you are acquiring the software on behalf of the Department of
              Defense, the software shall be classified as “Commercial Computer Software” and the
              Government shall have only “Restricted Rights” as defined in Clause 252.227-7013 (c) (1)
              of DFARs. Notwithstanding the foregoing, the authors grant the U.S. Government and
              others acting in its behalf permission to use and distribute the software in accordance with
              the terms specified in this license.
\end{lstlisting}

\subsection{This file is subject to the following copyright notice, which is different from the notice used}

\begin{lstlisting}[basicstyle=\ttfamily\scriptsize,breaklines=true]
This file is subject to the following copyright notice, which is different from the notice used
              elsewhere in Tcl, but roughly equivalent in meaning.
              Copyright (c) 1992,1993,1995,1996, Jens-Uwe Mager, Helios Software GmbH
              Not derived from licensed software.
              Permission is granted to freely use, copy, modify, and redistribute this software, provided
              that the author is not construed to be liable for any results of using the software,
              alterations are clearly marked as such, and this notice is not modified.
\end{lstlisting}

\subsection{Copyright (c) 1987-1994 The Regents of the University of California.}

\begin{lstlisting}[basicstyle=\ttfamily\scriptsize,breaklines=true]
Copyright (c) 1987-1994 The Regents of the University of California.
              Copyright (c) 1994-1995 Sun Microsystems, Inc.
              This software is copyrighted by the Regents of the University of California, Sun
              Microsystems, Inc., and other parties. The following terms apply to all files associated
              with the software unless explicitly disclaimed in individual files.
              The authors hereby grant permission to use, copy, modify, distribute, and license this
              software and its documentation for any purpose, provided that existing copyright notices
              are retained in all copies and that this notice is included verbatim in any distributions. No
              written agreement, license, or royalty fee is required for any of the authorized uses.
              Modifications to this software may be copyrighted by their authors and need not follow the
              licensing terms described here, provided that the new terms are clearly indicated on the
              first page of each file where they apply.
              IN NO EVENT SHALL THE AUTHORS OR DISTRIBUTORS BE LIABLE TO ANY PARTY
              FOR DIRECT, INDIRECT, SPECIAL, INCIDENTAL, OR CONSEQUENTIAL DAMAGES
              ARISING OUT OF THE USE OF THIS SOFTWARE, ITS DOCUMENTATION, OR ANY
              DERIVATIVES THEREOF, EVEN IF THE AUTHORS HAVE BEEN ADVISED OF THE
              POSSIBILITY OF SUCH DAMAGE.
              THE AUTHORS AND DISTRIBUTORS SPECIFICALLY DISCLAIM ANY WARRANTIES,
              INCLUDING, BUT NOT LIMITED TO, THE IMPLIED WARRANTIES OF
              MERCHANTABILITY, FITNESS FOR A PARTICULAR PURPOSE, AND NON-
              INFRINGEMENT. THIS SOFTWARE IS PROVIDED ON AN “AS IS” BASIS, AND THE
              AUTHORS AND DISTRIBUTORS HAVE NO OBLIGATION TO PROVIDE
              MAINTENANCE, SUPPORT, UPDATES, ENHANCEMENTS, OR MODIFICATIONS.
              RESTRICTED RIGHTS: Use, duplication or disclosure by the government is subject to
              the restrictions as set forth in subparagraph (c) (1) (ii) of the Rights in Technical Data and
              Computer Software Clause as DFARS 252.227-7013 and FAR 52.227-19.
ANSYS, Inc.

APPENDIX F — Third-Party Software Licenses                                   RedHawk User Manual | 979
Permission to use, copy, modify, and distribute this software and its documentation for
              any purpose and without fee is hereby granted, provided that the above copyright notice
              appear in all copies and that both that the copyright notice and warranty disclaimer
              appear in supporting documentation, and that the names of Lucent Technologies any of
              their entities not be used in advertising or publicity pertaining to distribution of the
              software without specific, written prior permission.
              Lucent disclaims all warranties with regard to this software, including all implied
              warranties of merchantability and fitness. In no event shall Lucent be liable for any
              special, indirect or consequential damages or any damages whatsoever resulting from
              loss of use, data or profits, whether in an action of contract, negligence or other tortuous
              action, arising out of or in connection with the use or performance of this software.
Gifsicle is copyright (C) 1997-2017 Eddie Kohler and licensed under GNU General Public
              License (GPL) Version 2. Please contact open.source@ansys.com for a copy of the
              Gifsicle source code.

              G Wave is copyright (C) Steve Tell and licensed under GNU General Public License
              (GPL) Version 2. Please contact open.source@ansys.com for a copy of the GNU
              Readline source code.

              GNU Readline is copyright (C) Free Software Foundation and licensed under GNU
              General Public License (GPL) Version 2. Please contact open.source@ansys.com for a
              copy of the GNU Readline source code.
              GNU GENERAL PUBLIC LICENSE
              Version 2, June 1991
              Copyright (C) 1989, 1991 Free Software Foundation, Inc.
              51 Franklin Street, Fifth Floor, Boston, MA 02110-1301, USA
Everyone is permitted to copy and distribute verbatim copies of this license document, but
              changing it is not allowed.
              Preamble
              The licenses for most software are designed to take away your freedom to share and
              change it. By contrast, the GNU General Public License is intended to guarantee your
              freedom to share and change free software--to make sure the software is free for all its
              users. This General Public License applies to most of the Free Software Foundation's
              software and to any other program whose authors commit to using it. (Some other Free
              Software Foundation software is covered by the GNU Lesser General Public License
              instead.) You can apply it to your programs, too.
              When we speak of free software, we are referring to freedom, not price. Our General
              Public Licenses are designed to make sure that you have the freedom to distribute copies
              of free software (and charge for this service if you wish), that you receive source code or
              can get it if you want it, that you can change the software or use pieces of it in new free
              programs; and that you know you can do these things.
              To protect your rights, we need to make restrictions that forbid anyone to deny you these
              rights or to ask you to surrender the rights. These restrictions translate to certain
              responsibilities for you if you distribute copies of the software, or if you modify it.
              For example, if you distribute copies of such a program, whether gratis or for a fee, you
              must give the recipients all the rights that you have. You must make sure that they, too,
APPENDIX F — Third-Party Software Licenses                                    RedHawk User Manual | 980
receive or can get the source code. And you must show them these terms so they know
              their rights.
              We protect your rights with two steps: (1) copyright the software, and (2) offer you this
              license which gives you legal permission to copy, distribute and/or modify the software.
              Also, for each author's protection and ours, we want to make certain that everyone
              understands that there is no warranty for this free software. If the software is modified by
              someone else and passed on, we want its recipients to know that what they have is not
              the original, so that any problems introduced by others will not reflect on the original
              authors' reputations.
              Finally, any free program is threatened constantly by software patents. We wish to avoid
              the danger that redistributors of a free program will individually obtain patent licenses, in
              effect making the program proprietary. To prevent this, we have made it clear that any
              patent must be licensed for everyone's free use or not licensed at all.
              The precise terms and conditions for copying, distribution and modification follow.
TERMS AND CONDITIONS FOR COPYING, DISTRIBUTION AND MODIFICATION
              0. This License applies to any program or other work which contains a notice placed by
              the copyright holder saying it may be distributed under the terms of this General Public
              License. The “Program”, below, refers to any such program or work, and a “work based
              on the Program” means either the Program or any derivative work under copyright law:
              that is to say, a work containing the Program or a portion of it, either verbatim or with
              modifications and/or translated into another language. (Hereinafter, translation is included
              without limitation in the term “modification”.) Each licensee is addressed as “you”.
              Activities other than copying, distribution and modification are not covered by this
              License; they are outside its scope. The act of running the Program is not restricted, and
              the output from the Program is covered only if its contents constitute a work based on the
              Program (independent of having been made by running the Program). Whether that is
              true depends on what the Program does.
              1. You may copy and distribute verbatim copies of the Program's source code as you
              receive it, in any medium, provided that you conspicuously and appropriately publish on
              each copy an appropriate copyright notice and disclaimer of warranty; keep intact all the
              notices that refer to this License and to the absence of any warranty; and give any other
              recipients of the Program a copy of this License along with the Program.
              You may charge a fee for the physical act of transferring a copy, and you may at your
              option offer warranty protection in exchange for a fee.
              2. You may modify your copy or copies of the Program or any portion of it, thus forming a
              work based on the Program, and copy and distribute such modifications or work under the
              terms of Section 1 above, provided that you also meet all of these conditions:
                      a) You must cause the modified files to carry prominent notices stating that you
                          changed the files and the date of any change.
                      b) You must cause any work that you distribute or publish, that in whole or in part
                          contains or is derived from the Program or any part thereof, to be licensed as
                          a whole at no charge to all third parties under the terms of this License.
                      c) If the modified program normally reads commands interactively when run, you
                            must cause it, when started running for such interactive use in the most
                            ordinary way, to print or display an announcement including an appropriate
                            copyright notice and a notice that there is no warranty (or else, saying that
                            you provide a warranty) and that users may redistribute the program under
                            these conditions, and telling the user how to view a copy of this License.
ANSYS, Inc.

APPENDIX F — Third-Party Software Licenses                                      RedHawk User Manual | 981
(Exception: if the Program itself is interactive but does not normally print such
                           an announcement, your work based on the Program is not required to print
                           an announcement.)
              These requirements apply to the modified work as a whole. If identifiable sections of that
              work are not derived from the Program, and can be reasonably considered independent
              and separate works in themselves, then this License, and its terms, do not apply to those
              sections when you distribute them as separate works. But when you distribute the same
              sections as part of a whole which is a work based on the Program, the distribution of the
              whole must be on the terms of this License, whose permissions for other licensees extend
              to the entire whole, and thus to each and every part regardless of who wrote it.
              Thus, it is not the intent of this section to claim rights or contest your rights to work written
              entirely by you; rather, the intent is to exercise the right to control the distribution of
              derivative or collective works based on the Program.
              In addition, mere aggregation of another work not based on the Program with the
              Program (or with a work based on the Program) on a volume of a storage or distribution
              medium does not bring the other work under the scope of this License.
              3. You may copy and distribute the Program (or a work based on it, under Section 2) in
              object code or executable form under the terms of Sections 1 and 2 above, provided that
              you also do one of the following:
                      a) Accompany it with the complete corresponding machine-readable source code,
                          which must be distributed under the terms of Sections 1 and 2 above on a
                          medium customarily used for software interchange; or,
                      b) Accompany it with a written offer, valid for at least three years, to give any third
                          party, for a charge no more than your cost of physically performing source
                          distribution, a complete machine-readable copy of the corresponding source
                          code, to be distributed under the terms of Sections 1 and 2 above on a
                          medium customarily used for software interchange; or,
                      c) Accompany it with the information you received as to the offer to distribute
                          corresponding source code. (This alternative is allowed only for
                          noncommercial distribution and only if you received the program in object
                          code or executable form with such an offer, in accord with Subsection b
                          above.)
              The source code for a work means the preferred form of the work for making
              modifications to it. For an executable work, complete source code means all the source
              code for all modules it contains, plus any associated interface definition files, plus the
              scripts used to control compilation and installation of the executable. However, as a
              special exception, the source code distributed need not include anything that is normally
              distributed (in either source or binary form) with the major components (compiler, kernel,
              and so on) of the operating system on which the executable runs, unless that component
              itself accompanies the executable.
              If distribution of executable or object code is made by offering access to copy from a
              designated place, then offering equivalent access to copy the source code from the same
              place counts as distribution of the source code, even though third parties are not
              compelled to copy the source along with the object code.
              4. You may not copy, modify, sub license, or distribute the Program except as expressly
              provided under this License. Any attempt otherwise to copy, modify, sub license or
              distribute the Program is void, and will automatically terminate your rights under this
              License. However, parties who have received copies, or rights, from you under this
              License will not have their licenses terminated so long as such parties remain in full
              compliance.
              5. You are not required to accept this License, since you have not signed it. However,
              nothing else grants you permission to modify or distribute the Program or its derivative
APPENDIX F — Third-Party Software Licenses                                    RedHawk User Manual | 982
works. These actions are prohibited by law if you do not accept this License. Therefore,
              by modifying or distributing the Program (or any work based on the Program), you
              indicate your acceptance of this License to do so, and all its terms and conditions for
              copying, distributing or modifying the Program or works based on it.
              6. Each time you redistribute the Program (or any work based on the Program), the
              recipient automatically receives a license from the original licensor to copy, distribute or
              modify the Program subject to these terms and conditions. You may not impose any
              further restrictions on the recipients' exercise of the rights granted herein. You are not
              responsible for enforcing compliance by third parties to this License.
              7. If, as a consequence of a court judgment or allegation of patent infringement or for any
              other reason (not limited to patent issues), conditions are imposed on you (whether by
              court order, agreement or otherwise) that contradict the conditions of this License, they do
              not excuse you from the conditions of this License. If you cannot distribute so as to satisfy
              simultaneously your obligations under this License and any other pertinent obligations,
              then as a consequence you may not distribute the Program at all. For example, if a patent
              license would not permit royalty-free redistribution of the Program by all those who
              receive copies directly or indirectly through you, then the only way you could satisfy both it
              and this License would be to refrain entirely from distribution of the Program.
              If any portion of this section is held invalid or unenforceable under any particular
              circumstance, the balance of the section is intended to apply and the section as a whole is
              intended to apply in other circumstances.
              It is not the purpose of this section to induce you to infringe any patents or other property
              right claims or to contest validity of any such claims; this section has the sole purpose of
              protecting the integrity of the free software distribution system, which is implemented by
              public license practices. Many people have made generous contributions to the wide
              range of software distributed through that system in reliance on consistent application of
              that system; it is up to the author/donor to decide if he or she is willing to distribute
              software through any other system and a licensee cannot impose that choice.
              This section is intended to make thoroughly clear what is believed to be a consequence of
              the rest of this License.
              8. If the distribution and/or use of the Program is restricted in certain countries either by
              patents or by copyrighted interfaces, the original copyright holder who places the Program
              under this License may add an explicit geographical distribution limitation excluding those
              countries, so that distribution is permitted only in or among countries not thus excluded. In
              such case, this License incorporates the limitation as if written in the body of this License.
              9. The Free Software Foundation may publish revised and/or new versions of the General
              Public License from time to time. Such new versions will be similar in spirit to the present
              version, but may differ in detail to address new problems or concerns.
              Each version is given a distinguishing version number. If the Program specifies a version
              number of this License which applies to it and “any later version”, you have the option of
              following the terms and conditions either of that version or of any later version published
              by the Free Software Foundation. If the Program does not specify a version number of
              this License, you may choose any version ever published by the Free Software
              Foundation.
              10. If you wish to incorporate parts of the Program into other free programs whose
              distribution conditions are different, write to the author to ask for permission. For software
              which is copyrighted by the Free Software Foundation, write to the Free Software
              Foundation; we sometimes make exceptions for this. Our decision will be guided by the
              two goals of preserving the free status of all derivatives of our free software and of
              promoting the sharing and reuse of software generally.
              NO WARRANTY
ANSYS, Inc.

APPENDIX F — Third-Party Software Licenses                                    RedHawk User Manual | 983
11. BECAUSE THE PROGRAM IS LICENSED FREE OF CHARGE, THERE IS NO
              WARRANTY FOR THE PROGRAM, TO THE EXTENT PERMITTED BY APPLICABLE
              LAW. EXCEPT WHEN OTHERWISE STATED IN WRITING THE COPYRIGHT
              HOLDERS AND/OR OTHER PARTIES PROVIDE THE PROGRAM “AS IS” WITHOUT
              WARRANTY OF ANY KIND, EITHER EXPRESSED OR IMPLIED, INCLUDING, BUT
              NOT LIMITED TO, THE IMPLIED WARRANTIES OF MERCHANTABILITY AND
              FITNESS FOR A PARTICULAR PURPOSE. THE ENTIRE RISK AS TO THE QUALITY
              AND PERFORMANCE OF THE PROGRAM IS WITH YOU. SHOULD THE PROGRAM
              PROVE DEFECTIVE, YOU ASSUME THE COST OF ALL NECESSARY SERVICING,
              REPAIR OR CORRECTION.
              12. IN NO EVENT UNLESS REQUIRED BY APPLICABLE LAW OR AGREED TO IN
              WRITING WILL ANY COPYRIGHT HOLDER, OR ANY OTHER PARTY WHO MAY
              MODIFY AND/OR REDISTRIBUTE THE PROGRAM AS PERMITTED ABOVE, BE
              LIABLE TO YOU FOR DAMAGES, INCLUDING ANY GENERAL, SPECIAL,
              INCIDENTAL OR CONSEQUENTIAL DAMAGES ARISING OUT OF THE USE OR
              INABILITY TO USE THE PROGRAM (INCLUDING BUT NOT LIMITED TO LOSS OF
              DATA OR DATA BEING RENDERED INACCURATE OR LOSSES SUSTAINED BY YOU
              OR THIRD PARTIES OR A FAILURE OF THE PROGRAM TO OPERATE WITH ANY
              OTHER PROGRAMS), EVEN IF SUCH HOLDER OR OTHER PARTY HAS BEEN
              ADVISED OF THE POSSIBILITY OF SUCH DAMAGES.
              END OF TERMS AND CONDITIONS
\end{lstlisting}

\subsection{CSPARSE is Copyright (C) 2006, Timothy Davis and licensed under GNU Library}

\begin{lstlisting}[basicstyle=\ttfamily\scriptsize,breaklines=true]
CSPARSE is Copyright (C) 2006, Timothy Davis and licensed under GNU Library
              General Public License (LGPL) Version 2.1. Please contact open.source@ansys.com for
              a copy of the CSPARSE source code.

              GDK, GTK+, and GLib are Copyright (C) The GNOME Project and licensed under GNU
              Library General Public License (LGPL) Version 2.1. Please contact
              open.source@ansys.com for a copy of the GDK or GLib source code.

              Libart is Copyright (C) Raph Levien and licensed under GNU Library General Public
              License (LGPL) Version 2.1. Please contact open.source@ansys.com for a copy of the
              Libart source code.

              The Qt GUI Toolkit is Copyright (C) 2015 The Qt Company Ltd. and licensed under GNU
              Library General Public License (LGPL) Version 2.1. Please contact
              open.source@ansys.com for a copy of the Qt GUI Toolkit source code.
\end{lstlisting}

\subsection{GNU LESSER GENERAL PUBLIC LICENSE}

\begin{lstlisting}[basicstyle=\ttfamily\scriptsize,breaklines=true]
GNU LESSER GENERAL PUBLIC LICENSE
              Version 2.1, February 1999
              Copyright (C) 1991, 1999 Free Software Foundation, Inc.
              51 Franklin St, Fifth Floor, Boston, MA 02110-1301 USA
              Everyone is permitted to copy and distribute verbatim copies of this license document, but
              changing it is not allowed.
              [This is the first released version of the Lesser GPL. It also counts as the successor of the
              GNU Library Public License, version 2, hence the version number 2.1.]
              Preamble
ANSYS, Inc.

APPENDIX F — Third-Party Software Licenses                                      RedHawk User Manual | 984
The licenses for most software are designed to take away your freedom to share and
              change it. By contrast, the GNU General Public Licenses are intended to guarantee your
              freedom to share and change free software--to make sure the software is free for all its
              users.
              This license, the Lesser General Public License, applies to some specially designated
              software packages--typically libraries--of the Free Software Foundation and other authors
              who decide to use it. You can use it too, but we suggest you first think carefully about
              whether this license or the ordinary General Public License is the better strategy to use in
              any particular case, based on the explanations below.
              When we speak of free software, we are referring to freedom of use, not price. Our
              General Public Licenses are designed to make sure that you have the freedom to
              distribute copies of free software (and charge for this service if you wish); that you receive
              source code or can get it if you want it; that you can change the software and use pieces
              of it in new free programs; and that you are informed that you can do these things.
              To protect your rights, we need to make restrictions that forbid distributors to deny you
              these rights or to ask you to surrender these rights. These restrictions translate to certain
              responsibilities for you if you distribute copies of the library or if you modify it.
              For example, if you distribute copies of the library, whether gratis or for a fee, you must
              give the recipients all the rights that we gave you. You must make sure that they, too,
              receive or can get the source code. If you link other code with the library, you must
              provide complete object files to the recipients, so that they can relink them with the library
              after making changes to the library and recompiling it. And you must show them these
              terms so they know their rights.
              We protect your rights with a two-step method: (1) we copyright the library, and (2) we
              offer you this license, which gives you legal permission to copy, distribute and/or modify
              the library.
              To protect each distributor, we want to make it very clear that there is no warranty for the
              free library. Also, if the library is modified by someone else and passed on, the recipients
              should know that what they have is not the original version, so that the original author's
              reputation will not be affected by problems that might be introduced by others.
              Finally, software patents pose a constant threat to the existence of any free program. We
              wish to make sure that a company cannot effectively restrict the users of a free program
              by obtaining a restrictive license from a patent holder. Therefore, we insist that any patent
              license obtained for a version of the library must be consistent with the full freedom of use
              specified in this license.
              Most GNU software, including some libraries, is covered by the ordinary GNU General
              Public License. This license, the GNU Lesser General Public License, applies to certain
              designated libraries, and is quite different from the ordinary General Public License. We
              use this license for certain libraries in order to permit linking those libraries into non-free
              programs.
              When a program is linked with a library, whether statically or using a shared library, the
              combination of the two is legally speaking a combined work, a derivative of the original
              library. The ordinary General Public License therefore permits such linking only if the
              entire combination fits its criteria of freedom. The Lesser General Public License permits
              more lax criteria for linking other code with the library.
              We call this license the “Lesser” General Public License because it does Less to protect
              the user's freedom than the ordinary General Public License. It also provides other free
              software developers Less of an advantage over competing non-free programs. These
              disadvantages are the reason we use the ordinary General Public License for many
APPENDIX F — Third-Party Software Licenses                                    RedHawk User Manual | 985
libraries. However, the Lesser license provides advantages in certain special
              circumstances.
              For example, on rare occasions, there may be a special need to encourage the widest
              possible use of a certain library, so that it becomes a de-facto standard. To achieve this,
              non-free programs must be allowed to use the library. A more frequent case is that a free
              library does the same job as widely used non-free libraries. In this case, there is little to
              gain by limiting the free library to free software only, so we use the Lesser General Public
              License.
              In other cases, permission to use a particular library in non-free programs enables a
              greater number of people to use a large body of free software. For example, permission to
              use the GNU C Library in non-free programs enables many more people to use the whole
              GNU operating system, as well as its variant, the GNU/Linux operating system.
              Although the Lesser General Public License is Less protective of the users' freedom, it
              does ensure that the user of a program that is linked with the Library has the freedom and
              the wherewithal to run that program using a modified version of the Library.
              The precise terms and conditions for copying, distribution and modification follow. Pay
              close attention to the difference between a “work based on the library” and a “work that
              uses the library”. The former contains code derived from the library, whereas the latter
              must be combined with the library in order to run.
              GNU LESSER GENERAL PUBLIC LICENSE TERMS AND CONDITIONS FOR
              COPYING, DISTRIBUTION AND MODIFICATION
              0. This License Agreement applies to any software library or other program which
              contains a notice placed by the copyright holder or other authorized party saying it may be
              distributed under the terms of this Lesser General Public License (also called “this
              License”). Each licensee is addressed as “you”.
              A “library” means a collection of software functions and/or data prepared so as to be
              conveniently linked with application programs (which use some of those functions and
              data) to form executables.
              The “Library”, below, refers to any such software library or work which has been
              distributed under these terms. A “work based on the Library” means either the Library or
              any derivative work under copyright law: that is to say, a work containing the Library or a
              portion of it, either verbatim or with modifications and/or translated straightforwardly into
              another language. (Hereinafter, translation is included without limitation in the term
              “modification”.)
              “Source code” for a work means the preferred form of the work for making modifications
              to it. For a library, complete source code means all the source code for all modules it
              contains, plus any associated interface definition files, plus the scripts used to control
              compilation and installation of the library.
              Activities other than copying, distribution and modification are not covered by this
              License; they are outside its scope. The act of running a program using the Library is not
              restricted, and output from such a program is covered only if its contents constitute a work
              based on the Library (independent of the use of the Library in a tool for writing it). Whether
              that is true depends on what the Library does and what the program that uses the Library
              does.
              1. You may copy and distribute verbatim copies of the Library's complete source code as
              you receive it, in any medium, provided that you conspicuously and appropriately publish
              on each copy an appropriate copyright notice and disclaimer of warranty; keep intact all
              the notices that refer to this License and to the absence of any warranty; and distribute a
              copy of this License along with the Library.You may charge a fee for the physical act of
              transferring a copy, and you may at your option offer warranty protection in exchange for a
              fee.
APPENDIX F — Third-Party Software Licenses                                      RedHawk User Manual | 986
2. You may modify your copy or copies of the Library or any portion of it, thus forming a
              work based on the Library, and copy and distribute such modifications or work under the
              terms of Section 1 above, provided that you also meet all of these conditions:
                      a) The modified work must itself be a software library.
                      b) You must cause the files modified to carry prominent notices stating that you
                          changed the files and the date of any change.
                      c) You must cause the whole of the work to be licensed at no charge to all third
                           parties under the terms of this License.
                      d) If a facility in the modified Library refers to a function or a table of data to be
                            supplied by an application program that uses the facility, other than as an
                            argument passed when the facility is invoked, then you must make a good
                            faith effort to ensure that, in the event an application does not supply such
                            function or table, the facility still operates, and performs whatever part of its
                            purpose remains meaningful. (For example, a function in a library to compute
                            square roots has a purpose that is entirely well-defined independent of the
                            application. Therefore, Subsection 2d requires that any application-supplied
                            function or table used by this function must be optional: if the application
                            does not supply it, the square root function must still compute square roots.)
              These requirements apply to the modified work as a whole. If identifiable sections of that
              work are not derived from the Library, and can be reasonably considered independent
              and separate works in themselves, then this License, and its terms, do not apply to those
              sections when you distribute them as separate works. But when you distribute the same
              sections as part of a whole which is a work based on the Library, the distribution of the
              whole must be on the terms of this License, whose permissions for other licensees extend
              to the entire whole, and thus to each and every part regardless of who wrote it.
              Thus, it is not the intent of this section to claim rights or contest your rights to work written
              entirely by you; rather, the intent is to exercise the right to control the distribution of
              derivative or collective works based on the Library.
              In addition, mere aggregation of another work not based on the Library with the Library (or
              with a work based on the Library) on a volume of a storage or distribution medium does
              not bring the other work under the scope of this License.
              3. You may opt to apply the terms of the ordinary GNU General Public License instead of
              this License to a given copy of the Library. To do this, you must alter all the notices that
              refer to this License, so that they refer to the ordinary GNU General Public License,
              version 2, instead of to this License. (If a newer version than version 2 of the ordinary
              GNU General Public License has appeared, then you can specify that version instead if
              you wish.) Do not make any other change in these notices. Once this change is made in a
              given copy, it is irreversible for that copy, so the ordinary GNU General Public License
              applies to all subsequent copies and derivative works made from that copy. This option is
              useful when you wish to copy part of the code of the Library into a program that is not a
              library.
              4. You may copy and distribute the Library (or a portion or derivative of it, under Section
              2) in object code or executable form under the terms of Sections 1 and 2 above provided
              that you accompany it with the complete corresponding machine-readable source code,
              which must be distributed under the terms of Sections 1 and 2 above on a medium
              customarily used for software interchange.
              If distribution of object code is made by offering access to copy from a designated place,
              then offering equivalent access to copy the source code from the same place satisfies the
              requirement to distribute the source code, even though third parties are not compelled to
              copy the source along with the object code.
ANSYS, Inc.

APPENDIX F — Third-Party Software Licenses                                        RedHawk User Manual | 987
5. A program that contains no derivative of any portion of the Library, but is designed to
              work with the Library by being compiled or linked with it, is called a “work that uses the
              Library”. Such a work, in isolation, is not a derivative work of the Library, and therefore
              falls outside the scope of this License.
              However, linking a “work that uses the Library” with the Library creates an executable that
              is a derivative of the Library (because it contains portions of the Library), rather than a
              “work that uses the library”. The executable is therefore covered by this License.
              Section 6 states terms for distribution of such executables.
              When a “work that uses the Library” uses material from a header file that is part of the
              Library, the object code for the work may be a derivative work of the Library even though
              the source code is not. Whether this is true is especially significant if the work can be
              linked without the Library, or if the work is itself a library. The threshold for this to be true is
              not precisely defined by law. If such an object file uses only numerical parameters, data
              structure layouts and accessors, and small macros and small inline functions (ten lines or
              less in length), then the use of the object file is unrestricted, regardless of whether it is
              legally a derivative work. (Executables containing this object code plus portions of the
              Library will still fall under Section 6.) Otherwise, if the work is a derivative of the Library,
              you may distribute the object code for the work under the terms of Section 6. Any
              executables containing that work also fall under Section 6, whether or not they are linked
              directly with the Library itself.
              6. As an exception to the Sections above, you may also combine or link a “work that uses
              the Library” with the Library to produce a work containing portions of the Library, and
              distribute that work under terms of your choice, provided that the terms permit
              modification of the work for the customer's own use and reverse engineering for
              debugging such modifications.
              You must give prominent notice with each copy of the work that the Library is used in it
              and that the Library and its use are covered by this License. You must supply a copy of
              this License. If the work during execution displays copyright notices, you must include the
              copyright notice for the Library among them, as well as a reference directing the user to
              the copy of this License. Also, you must do one of these things:
                      a) Accompany the work with the complete corresponding machine-readable
                          source code for the Library including whatever changes were used in the
                          work (which must be distributed under Sections 1 and 2 above); and, if the
                          work is an executable linked with the Library, with the complete machine-
                          readable “work that uses the Library”, as object code and/or source code, so
                          that the user can modify the Library and then relink to produce a modified
                          executable containing the modified Library. (It is understood that the user
                          who changes the contents of definitions files in the Library will not necessarily
                          be able to recompile the application to use the modified definitions.)
                      b) Use a suitable shared library mechanism for linking with the Library. A suitable
                          mechanism is one that (1) uses at run time a copy of the library already
                          present on the user's computer system, rather than copying library functions
                          into the executable, and (2) will operate properly with a modified version of
                          the library, if the user installs one, as long as the modified version is
                          interface-compatible with the version that the work was made with.
                         c) Accompany the work with a written offer, valid for at least three years, to give
                            the same user the materials specified in Subsection 6a, above, for a charge
                            no more than the cost of performing this distribution.
                      d) If distribution of the work is made by offering access to copy from a designated
                            place, offer equivalent access to copy the above specified materials from the
                            same place.
ANSYS, Inc.

APPENDIX F — Third-Party Software Licenses                                    RedHawk User Manual | 988
e) Verify that the user has already received a copy of these materials or that you
                          have already sent this user a copy.
              For an executable, the required form of the “work that uses the Library” must include any
              data and utility programs needed for reproducing the executable from it. However, as a
              special exception, the materials to be distributed need not include anything that is
              normally distributed (in either source or binary form) with the major components (compiler,
              kernel, and so on) of the operating system on which the executable runs, unless that
              component itself accompanies the executable.
              It may happen that this requirement contradicts the license restrictions of other proprietary
              libraries that do not normally accompany the operating system. Such a contradiction
              means you cannot use both them and the Library together in an executable that you
              distribute.
              7. You may place library facilities that are a work based on the Library side-by-side in a
              single library together with other library facilities not covered by this License, and
              distribute such a combined library, provided that the separate distribution of the work
              based on the Library and of the other library facilities is otherwise permitted, and provided
              that you do these two things:
                      a) Accompany the combined library with a copy of the same work based on the
                          Library, uncombined with any other library facilities. This must be distributed
                          under the terms of the Sections above.
                      b) Give prominent notice with the combined library of the fact that part of it is a
                          work based on the Library, and explaining where to find the accompanying
                          uncombined form of the same work.
              8. You may not copy, modify, sub license, link with, or distribute the Library except as
              expressly provided under this License. Any attempt otherwise to copy, modify, sub
              license, link with, or distribute the Library is void, and will automatically terminate your
              rights under this License. However, parties who have received copies, or rights, from you
              under this License will not have their licenses terminated so long as such parties remain
              in full compliance.
              9. You are not required to accept this License, since you have not signed it. However,
              nothing else grants you permission to modify or distribute the Library or its derivative
              works. These actions are prohibited by law if you do not accept this License. Therefore,
              by modifying or distributing the Library (or any work based on the Library), you indicate
              your acceptance of this License to do so, and all its terms and conditions for copying,
              distributing or modifying the Library or works based on it.
              10. Each time you redistribute the Library (or any work based on the Library), the recipient
              automatically receives a license from the original licensor to copy, distribute, link with or
              modify the Library subject to these terms and conditions. You may not impose any further
              restrictions on the recipients' exercise of the rights granted herein. You are not
              responsible for enforcing compliance by third parties with this License.
              11. If, as a consequence of a court judgment or allegation of patent infringement or for any
              other reason (not limited to patent issues), conditions are imposed on you (whether by
              court order, agreement or otherwise) that contradict the conditions of this License, they do
              not excuse you from the conditions of this License. If you cannot distribute so as to satisfy
              simultaneously your obligations under this License and any other pertinent obligations,
              then as a consequence you may not distribute the Library at all. For example, if a patent
              license would not permit royalty-free redistribution of the Library by all those who receive
              copies directly or indirectly through you, then the only way you could satisfy both it and
              this License would be to refrain entirely from distribution of the Library.
              If any portion of this section is held invalid or unenforceable under any particular
              circumstance, the balance of the section is intended to apply, and the section as a whole
              is intended to apply in other circumstances.
APPENDIX F — Third-Party Software Licenses                                    RedHawk User Manual | 989
It is not the purpose of this section to induce you to infringe any patents or other property
              right claims or to contest validity of any such claims; this section has the sole purpose of
              protecting the integrity of the free software distribution system which is implemented by
              public license practices. Many people have made generous contributions to the wide
              range of software distributed through that system in reliance on consistent application of
              that system; it is up to the author/donor to decide if he or she is willing to distribute
              software through any other system and a licensee cannot impose that choice.
              This section is intended to make thoroughly clear what is believed to be a consequence of
              the rest of this License.
              12. If the distribution and/or use of the Library is restricted in certain countries either by
              patents or by copyrighted interfaces, the original copyright holder who places the Library
              under this License may add an explicit geographical distribution limitation excluding those
              countries, so that distribution is permitted only in or among countries not thus excluded. In
              such case, this License incorporates the limitation as if written in the body of this License.
              13. The Free Software Foundation may publish revised and/or new versions of the Lesser
              General Public License from time to time. Such new versions will be similar in spirit to the
              present version, but may differ in detail to address new problems or concerns.
              Each version is given a distinguishing version number. If the Library specifies a version
              number of this License which applies to it and “any later version”, you have the option of
              following the terms and conditions either of that version or of any later version published
              by the Free Software Foundation. If the Library does not specify a license version number,
              you may choose any version ever published by the Free Software Foundation.
              14. If you wish to incorporate parts of the Library into other free programs whose
              distribution conditions are incompatible with these, write to the author to ask for
              permission. For software which is copyrighted by the Free Software Foundation, write to
              the Free Software Foundation; we sometimes make exceptions for this. Our decision will
              be guided by the two goals of preserving the free status of all derivatives of our free
              software and of promoting the sharing and reuse of software generally.
              NO WARRANTY
              15. BECAUSE THE LIBRARY IS LICENSED FREE OF CHARGE, THERE IS NO
              WARRANTY FOR THE LIBRARY, TO THE EXTENT PERMITTED BY APPLICABLE
              LAW. EXCEPT WHEN OTHERWISE STATED IN WRITING THE COPYRIGHT
              HOLDERS AND/OR OTHER PARTIES PROVIDE THE LIBRARY “AS IS” WITHOUT
              WARRANTY OF ANY KIND, EITHER EXPRESSED OR IMPLIED, INCLUDING, BUT
              NOT LIMITED TO, THE IMPLIED WARRANTIES OF MERCHANTABILITY AND
              FITNESS FOR A PARTICULAR PURPOSE. THE ENTIRE RISK AS TO THE QUALITY
              AND PERFORMANCE OF THE LIBRARY IS WITH YOU. SHOULD THE LIBRARY
              PROVE DEFECTIVE, YOU ASSUME THE COST OF ALL NECESSARY SERVICING,
              REPAIR OR CORRECTION.
              16. IN NO EVENT UNLESS REQUIRED BY APPLICABLE LAW OR AGREED TO IN
              WRITING WILL ANY COPYRIGHT HOLDER, OR ANY OTHER PARTY WHO MAY
              MODIFY AND/OR REDISTRIBUTE THE LIBRARY AS PERMITTED ABOVE, BE LIABLE
              TO YOU FOR DAMAGES, INCLUDING ANY GENERAL, SPECIAL, INCIDENTAL OR
              CONSEQUENTIAL DAMAGES ARISING OUT OF THE USE OR INABILITY TO USE
              THE LIBRARY (INCLUDING BUT NOT LIMITED TO LOSS OF DATA OR DATA BEING
              RENDERED INACCURATE OR LOSSES SUSTAINED BY YOU OR THIRD PARTIES
              OR A FAILURE OF THE LIBRARY TO OPERATE WITH ANY OTHER SOFTWARE),
              EVEN IF SUCH HOLDER OR OTHER PARTY HAS BEEN ADVISED OF THE
              POSSIBILITY OF SUCH DAMAGES.
              END OF TERMS AND CONDITIONS
\end{lstlisting}
